\chapter{The DAT File Format Specification}
\label{chap:datref}

The solver input data is contained in a flat text file (the ``DAT file,'' usually with extension \verb+.dat+) with its own format as documented here.

Notes on the format:
\begin{Snippet}
# DAT files are ASCII and whitespace-delimited with one record per line
# Labels (From, At, To) denote positions (Points or POS records)
# Each POS record (Point) must have a unique label;
#    but not all labels need have a POS record.
# keywords are in ALL CAPS
# POS ... CONS c[c'] constrain the position solution to a plane[line]
#    using c,c' = one of{XYZ} OR {NEU} and c != c'
# Note that covariance is required for POS with ADJ
# FIX : Point is constant (not adjusted);
# ADJ : use position as a priori information and adjust control point;
# EST or blank: use position only as a priori.
# LUNIT means linear unit MM|M|CM|KM|FT
# AUNIT means angular unit RAD|DEG|SOA but not DMS
# DMS denotes 'deg/min/sec' (int/int/float) and follows angle given by deg min sec
# NB in the angle records "angle sig AUNIT" means angle and sigma
#   both have unit AUNIT; consider "angle AUNIT sig AUNIT"
# fcsig acsig and tcsig are centering errors on From, At and To stations in LUNITs
# CORR is optional but must be followed by corrections in the order shown;
#   use zero placeholders.
# ht[f|t] is height at From|To; target height diff. is implemented using From+To
# refract is a float refraction, and CURV, OHC and GEOID are keywords meaning
#    apply curvature, orthometric height and (DOV, undulation) corrections,
#    respectively.
# Title is for output only
# comment lines begin with '#' and are ignored; also #-to-EOL is ignored
#
# Tags: each measurement record (all measurements+DirSet+Dir+POS ADJ)
#   can be given an optional tag with the field tag=TAG
#  where TAG is the user's tag. This field MUST BE THE LAST FIELD on the line.
#   This tag will be used in generating
#  "data names" in the solver, rather than simply numbering them;
#   the number will be used when a tag is not found.
#  The user MUST use unique tags; failure to do so will cause the solver
#   to throw an exception.
\end{Snippet}

The records are defined as follows, each a single line (even though some lines here are wrapped), beginning with a keyword.

Configuration records.
\begin{Snippet}
# Configuration: output precision; problem dimension (2D is XY);
#   convergence criteria
TITLE title # quotes optional
PREC  eps UNIT [eps UNIT] # output precision; linear and/or angular units
DIM 2|3  # dimension of the problem; 2D is XY only, i.e. Z's ignored
CONV [n ITER] [d CONV] # convergence criteria; either or both
CONFIDENCE alpha       # confidence for Chi squ test (0 < alpha < 1)
OUT [NOAPV] [APQUIT]   # output NOAPV = do not scale covariance with APV
OUT [NOAPV] [APQUIT]   # output APQUIT = quit after ComputeAPriori()
GEOIDFILE filename     # Filename for the gridded EGM08 geoid file
INTERPOLATION method   # Geoid Interpolation method (bicubic, bilinear)
HASH string            # GUI may want to pass hash through solver to binary file
COMMENT ...            # comment that is echoed in output file
GEOID <file>           # geoid file
\end{Snippet}

Position (site or Point) record.
\begin{Snippet}
# no corrections; constraints c one[two] of XYZ|NEU
POS label X Y Z [covxx xy xz yy yz zz] LUNIT [FIX|ADJ|EST] [CONS c[c]]
# undocumented
POS label D M S N|S D M S E|W Ht LUNIT [covxx xy xz yy yz zz [NEU] LUNIT]
         [FIX|ADJ|EST] [CONS c[c]]

POS label X Y Z covxx xy xz yy yz zz LUNIT FIX|ADJ|EST CONS c[c]
POS label X Y Z covxx xy xz yy yz zz LUNIT FIX|ADJ|EST
POS label X Y Z covxx xy xz yy yz zz LUNIT             CONS c[c]
POS label X Y Z LUNIT                      FIX|EST     CONS c[c]
POS label X Y Z LUNIT                      FIX|EST
POS label X Y Z LUNIT                                  CONS c[c]
POS label D M S N|S D M S E|W Ht LUNIT covxx xy xz yy yz zz NEU LUNIT FIX|ADJ|EST
         CONS c[c]
POS label D M S N|S D M S E|W Ht LUNIT covxx xy xz yy yz zz NEU LUNIT FIX|ADJ|EST
POS label D M S N|S D M S E|W Ht LUNIT covxx xy xz yy yz zz NEU LUNIT            
         CONS c[c]
POS label D M S N|S D M S E|W Ht LUNIT covxx xy xz yy yz zz     LUNIT FIX|ADJ|EST
         CONS c[c]
POS label D M S N|S D M S E|W Ht LUNIT covxx xy xz yy yz zz     LUNIT FIX|ADJ|EST
POS label D M S N|S D M S E|W Ht LUNIT covxx xy xz yy yz zz     LUNIT            
         CONS c[c]
POS label D M S N|S D M S E|W Ht LUNIT                                FIX|EST    
         CONS c[c]
POS label D M S N|S D M S E|W Ht LUNIT                                FIX|EST
POS label D M S N|S D M S E|W Ht LUNIT                                           
         CONS c[c]
\end{Snippet}

Measurement records (again, one per line even though some are wrapped here).
\begin{Snippet}
# 3-D delta XYZ (Delta)
DEL labFr labTo dX dY dZ covxx xy xz yy yz zz [asig [PPM]] [fcsig tcsig] LUNIT
        [CORR htf htt [LUNIT]]

# Distance (Dist)
DIS labFr labTo distance sig [asig [PPM]] [fcsig tcsig] LUNIT
        [CORR htf htt [LUNIT] [refract]]

# height (Height) GEOID if orthometric and undulation correction must be applied
HGT labFr labTo height sig [asig [PPM]] LUNIT [CORR [refract]
        [CURV] [GEOID]]

# Azimuth (Azimuth)
AZM labFr labTo angle [min sec DMS|AUNIT] sig [asig] AUNIT [fcsig tcsig LUNIT]
        [CORR htf htt [LUNIT] [GEOID]]

# Horizonal angle (HAngle)
HAN labFr labAt labTo angle [min sec DMS|AUNIT] sig [asig] AUNIT [fc ac tc LUNIT]
        [CORR htf htt [LUNIT] [GEOID]]

# Direction set and Direction (each complete set used to create set of HANs)
# Set = { One DIRSET + >1 DIR with one "group" string};
#   group for each set must be unique and identical throughout set
# Zero or one DIR in a set -> set is ignored - no HANs can be constructed
# DIRs where group does not appear in a DIRSET are ignored
DIRSET group labAt [ac LUNIT] [CORR GEOID]
DIR group labTo angle [min sec DMS|AUNIT] sig [asig] AUNIT [tc LUNIT]
        [CORR htt [LUNIT]]

# Vertical angle (VAngle)
VAN labFr labTo angle [min sec DMS|AUNIT] sig [asig] AUNIT [fc tc LUNIT]
        [CORR htf htt [LUNIT] [refract] [GEOID]]

# Zenith angle (ZAngle)
ZAN labFr labTo angle [min sec DMS|AUNIT] sig [asig] AUNIT [ac tc LUNIT]
        [CORR htf htt [LUNIT] [refract] [GEOID]]
\end{Snippet}

A table showing which records have centering, sigma, PPM and various corrections.
\begin{Snippet}
      centering    --sigs--   -------corrections---------  notes
      FC  AC  TC   sig  PPM   HI/HT DOV Undul Refrac Curv

POS    -   -   -    y    -      -    -    -     -     -    also constraints

DEL    y   -   y    y    y      y    -    -     -     -    GNSS relative position

DIS    y   -   y    y    y      y    -    -     y     -

HGT    -   -   -    y    y      -    y    y     y     y    Orthometric/ellipsoid
                                                             as [GEOID]

AZM    y   -   y    y    -      y    y    -     -     -

HAN    y   y   y    y    -      y    y    -     -     -

VAN    y   -   y    y    -      y    y    y     y     -    keyword GEOID:
                                                             geodetic or not

ZAN    y   -   y    y    -      y    y    y     y     -    keyword GEOID:
                                                             geodetic or not
\end{Snippet}

Examples.
\begin{Snippet}
# example
DIRSET A00 TOWER1 0.001 M
DIR A00 TOWER2  0  0   0.0 DMS 15 SOA 0.001 M      # no heights
DIR A00 TOWER3 62 41  38.2 DMS 15 SOA 0.001 M      # no ht
DIR A00 TOWER    95 57  16.0 DMS 15 SOA 0.001 M      # no ht

DIRSET A01 TOWER2 0.001 M
DIR A01 TOWER1   0  0  0.0 DMS 15 SOA 0.001 M CORR 1.508 M
DIR A01 TOWER3 311 46 17.9 DMS 15 SOA 0.001 M CORR 1.407 M
DIR A01 TOWER    314 29 32.2 DMS 15 SOA 0.001 M CORR 0.000 M

DIRSET A02 TOWER3 0.001 M
DIR A02 TOWER1   0  0  0.0 DMS 15 SOA 0.001 M CORR 1.508 M
DIR A02 TOWER2  69  4 34.5 DMS 15 SOA 0.001 M CORR 1.449 M
DIR A02 TOWER    255 47 15.4 DMS 15 SOA 0.001 M CORR 0.000 M
DIR A02 TOWERB    255 59 27.0 DMS 60 SOA 0.001 M CORR 0.000 M

# Extraction. block string creates groups, groups are extracted with
#    final nominal value and full covariance.
# Option to store solution in binary file, reload just for extraction.
# Does not apply to data, as such, b/c measurements of same thing can be repeated.
EXTR blk POS label (3)
EXTR blk DEL labFr labTo (3)
EXTR blk DIS labFr labTo
EXTR blk HGT labFr labTo
EXTR blk AZM labFr labTo
EXTR blk HAN labFr labAt labTo
EXTR blk VAN labFr labTo
EXTR blk ZAN labFr labTo
EXTR blk DIRSET group
\end{Snippet}
